% sections/problem.tex -- problem and wider cultural stakes (shared).
% B: "The problem and its wider stakes"; C: "Introduction".
% Claims and sources: source_verification.md, section E.
GenAI systems increasingly mediate access to and participation in some forms
of cultural production. Their training data, defaults, moderation rules,
access conditions and contractual terms are shaped largely by platform
companies \citep{crawford2021}. I therefore treat GenAI not as a neutral
tool or autonomous artist, but as part of a cultural infrastructure that
includes interfaces, institutions, facilitators and unequal conditions of
access.

Three mechanisms give this infrastructure social stakes. The first is \emph{platformized interpretive
authority}. \citet[276]{richards2025} argues that the ``safe and helpful''
assistant functions ``as a proxy for authority''; Caines found that
subscription tiers, privacy conditions and opacity shaped what she and her
students could make \citep[187--188, 194]{caines2026steaming}; and Caines
and Watson describe a project whose outcome lay ``completely in the hands of
the corporations that designed the software''
\citep[267]{caines2023consultation}. The second is \emph{cultural
legibility}: models may privilege, flatten, standardize, mistranslate,
misrecognize or pass over culturally situated expression. South Asian
participants found that text-to-image models reproduced an outsider's gaze
on their cultures \citep{qadri2023}, and AI writing suggestions moved Indian
writers toward Western styles \citep{agarwal2025}; each finding concerns
particular models and tasks. The third is the gap between
\emph{technical provenance and social authorship}: a system can record who
supplied what while collaborators credit someone else, and artists' work has
entered training data without their consent \citep{jiang2023}.
Environmental cost, labour displacement and misinformation form the wider
context \citep[182]{caines2026steaming} but are not objects of this study.

Participation and authorship matter because they determine who takes part
in meaning-making, whose contribution remains recognizable, who can revise
or refuse a model's transformation, who receives public credit and who is
held responsible when collaborative work circulates. They are where platform
power becomes practical in everyday cultural production. One qualitative
study cannot explain AI's effect on society, but situated research-creation
can make these processes observable and yield analytically transferable
concepts.
